\documentclass[twocolumn,11pt,a4paper]{article}

\usepackage[utf8]{inputenc}
\usepackage[T1]{fontenc}
\usepackage{lmodern}
\usepackage[a4paper,top=2cm,bottom=2.2cm,left=1.8cm,right=1.8cm,columnsep=0.5cm]{geometry}
\usepackage{amsmath,amssymb,amsthm,mathtools}
\usepackage{bm}
\usepackage{booktabs}
\usepackage{array}
\usepackage{enumitem}
\usepackage{graphicx}
\usepackage{xcolor}
\usepackage{microtype}
\usepackage[colorlinks=true,linkcolor=blue!60!black,citecolor=blue!60!black,urlcolor=blue!60!black]{hyperref}
\usepackage{cite}
\usepackage{url}
\usepackage{fancyhdr}
\usepackage{abstract}
\usepackage{stmaryrd}
\usepackage{algorithm}
\usepackage{algpseudocode}

\pagestyle{fancy}
\fancyhf{}
\fancyhead[L]{\small\textit{Dual-Domain Execution Model}}
\fancyhead[R]{\small\thepage}
\renewcommand{\headrulewidth}{0.4pt}

\newtheorem{theorem}{Theorem}[section]
\newtheorem{lemma}[theorem]{Lemma}
\newtheorem{corollary}[theorem]{Corollary}
\newtheorem{proposition}[theorem]{Proposition}
\newtheorem{definition}[theorem]{Definition}
\newtheorem{remark}[theorem]{Remark}
\newtheorem{example}[theorem]{Example}
\newtheorem{axiom}{Axiom}

\newcommand{\DP}{\Delta P}
\newcommand{\Sk}{S_k}
\newcommand{\St}{S_t}
\newcommand{\Se}{S_e}
\newcommand{\Stuple}{\mathbf{S}}
\newcommand{\calC}{\mathcal{C}}
\newcommand{\calD}{\mathcal{D}}
\newcommand{\calA}{\mathcal{A}}
\newcommand{\calT}{\mathcal{T}}
\newcommand{\calP}{\mathcal{P}}
\newcommand{\calS}{\mathcal{S}}
\newcommand{\calE}{\mathcal{E}}
\newcommand{\calV}{\mathcal{V}}
\newcommand{\calF}{\mathcal{F}}
\newcommand{\Tref}{T_{\mathrm{ref}}}
\newcommand{\trec}{t_{\mathrm{rec}}}
\newcommand{\fref}{f_{\mathrm{ref}}}
\newcommand{\Mcyc}{M}
\newcommand{\llb}{\llbracket}
\newcommand{\rrb}{\rrbracket}
\DeclareMathOperator{\cell}{cell}
\DeclareMathOperator{\mean}{mean}

\title{%
  \LARGE\bfseries Dual-Domain Execution Model:\\[6pt]
  S-Entropy Bidirectional Search in Temporal Programs
  with Virtual Substate Decomposition
}

\author{Anonymous}

\date{\today}

\begin{document}

\twocolumn[{%
\maketitle
\begin{@twocolumnfalse}

\begin{abstract}
We present a unified computational model combining temporal programming (timing-based execution with parser-free security), S-entropy state representation (3D bounded partition space), and virtual substate decomposition (off-shell mathematical components). The key innovation is a \emph{dual-domain execution architecture} where programs simultaneously execute in time-domain (via timing cells) and frequency-domain (via virtual substate components), mediated by a mean-recovery constraint. We introduce the \emph{S-Entropy Bidirectional Dijkstra (SEBD) algorithm}, which searches forward through time-domain subtasks and backward through frequency-domain virtual substates, meeting at S-entropy coordinates to find minimum-cost equivalence chains. The algorithm achieves three simultaneous objectives: (1) deterministic, provably-secure execution (no payload parser); (2) optimal path finding in 3D partition space; (3) parallel subtask execution with guaranteed coherence via mean-recovery. We prove that the algorithm terminates, maintains correctness under the mean-recovery constraint, and enables detection of ``impossible states'' (partition addresses with no physical realization). Applications include real-time control systems, distributed positioning, and privacy-preserving computation.
\end{abstract}

\vspace{1em}

\end{@twocolumnfalse}
}]

\section{Introduction}
\label{sec:intro}

Modern computing systems face a fundamental tension: expressive power (ability to handle rich data) versus security (immunity to injection attacks and parsing exploits). Rich message formats enable flexibility but create vast attack surfaces. Bare timing signals are secure but expressive only through narrow channels.

We resolve this tension by introducing a \emph{dual-domain execution model} that executes simultaneously in two complementary spaces: the \emph{temporal domain} (where programs are sequences of timing cells, parser-free and structurally secure) and the \emph{frequency domain} (where the same computation manifests as a spectrum of virtual substate components with off-shell degrees of freedom). The two domains are mathematically equivalent via Fourier duality, but more importantly, they are \emph{causally mediated} by a constraint we call \emph{mean recovery}: the mean of all virtual components in the frequency domain must equal the physical state observable in the temporal domain.

\subsection{Motivating Example}

Consider a distributed positioning system where each node measures timing deviations $\DP$ from a local oscillator. Traditional approaches:
\begin{itemize}[nosep]
  \item \textbf{Temporal only}: Classify $\DP$ into cells, dispatch actions. Simple but 1-dimensional.
  \item \textbf{Frequency only}: Fourier-transform the signal, solve in frequency space. Rich but loses timing guarantees.
\end{itemize}

Our approach: Represent the node's state as a 3D vector $\Stuple = (\Sk, \St, \Se) \in [0,1]^3$ (S-entropy coordinates), execute timing cells in the temporal view, and simultaneously access decompositions of $\Stuple$ into $N$ virtual components (which may lie outside $[0,1]^3$) in the frequency view. The physical state is always recoverable as the mean of all virtual components.

The shortest path from current state to goal state can then be found by:
\begin{enumerate}[nosep]
  \item Searching \emph{forward} in time-domain subtasks (Dijkstra from current $\Stuple$)
  \item Searching \emph{backward} from the goal in frequency-domain virtual substates
  \item Meeting at an intermediate S-entropy coordinate where the cost is minimized
\end{enumerate}

This bidirectional search reduces the effective search space by orders of magnitude compared to unidirectional approaches.

\subsection{Contributions}

This paper makes the following contributions:

\begin{enumerate}[leftmargin=1.5em,topsep=2pt,itemsep=1pt]
  \item \textbf{Dual-domain architecture.} We formalize a computational model where temporal programs (timing cells and actions) are the primary interface, and frequency-domain virtual substates are an equivalent dual view. Both are coordinated by mean-recovery.

  \item \textbf{S-entropy representation.} We show that a 3-tuple $(\Sk, \St, \Se) \in [0,1]^3$ (kinetic, thermal, energetic entropy components) provides a natural state space for bounded oscillatory systems, generalizing the 1D timing deviation.

  \item \textbf{Virtual substate framework.} We develop the formal structure of off-shell virtual decompositions, prove the Miracle Principle (any physical state is compatible with any set of local off-shell components), and establish the mean-recovery constraint.

  \item \textbf{S-Entropy Bidirectional Dijkstra (SEBD) algorithm.} We introduce a shortest-path algorithm that searches forward in time domain and backward in frequency domain, meeting at optimal S-entropy coordinates. We prove termination, correctness, and optimality.

  \item \textbf{Security analysis.} We show that the temporal-domain view is attack-resistant (no parser, deterministic dispatch), while the frequency-domain view provides rich computational structure.

  \item \textbf{Coherence guarantee.} We prove that mean-recovery enforces global coherence: no matter how virtual components are distributed off-shell, the physical observable remains bounded and consistent.
\end{enumerate}

\section{Oscillatory Systems and Timing Fundamentals}
\label{sec:oscillator}

\subsection{Reference Oscillators}

\begin{definition}[Reference Oscillator]
A \emph{reference oscillator} is a tuple $\mathcal{O} = (\fref, \epsilon, M_0)$ where:
\begin{itemize}[nosep,leftmargin=1em]
  \item $\fref \in \mathbb{R}_{>0}$ is the nominal frequency (Hz),
  \item $\epsilon \in [0,1)$ is the fractional frequency stability bound,
  \item $M_0 \in \mathbb{N}$ is the initial cycle count.
\end{itemize}
The reference times are $\Tref(k) = k/\fref$ for $k \geq M_0$, and the cycle count at time $t$ is $\Mcyc(t) = M_0 + \lfloor (t - t_0) \cdot \fref \rfloor$.
\end{definition}

\subsection{Timing Deviations}

\begin{definition}[Timing Deviation]
For a pulse received at time $\trec$, when the oscillator predicts arrival at $\Tref(k)$, the \emph{timing deviation} is:
\begin{equation}
  \DP(k) = \Tref(k) - \trec.
  \label{eq:dp-def}
\end{equation}
In practice, $\DP \in \calD = [-\Delta_{\max}, +\Delta_{\max}]$ where $\Delta_{\max} = (\epsilon + \delta)/\fref$ for jitter bound $\delta$.
\end{definition}

\begin{theorem}[Monotone Cycle Count]
For $t_1 < t_2$, $\Mcyc(t_1) \leq \Mcyc(t_2)$. Moreover, if a signal with $\DP_1$ is recorded at $\Mcyc(t_1)$ and replayed at $\Mcyc(t_2) > \Mcyc(t_1)$, the new deviation is $\DP_2 = \DP_1 + (\Mcyc(t_2) - \Mcyc(t_1))/\fref \neq \DP_1$.
\end{theorem}

This monotonicity guarantees structural replay resistance: any replayed signal appears shifted in $\calD$ and falls in a different cell.

\section{S-Entropy State Representation}
\label{sec:sentropy}

\subsection{Three-Dimensional Partition State}

The 1D timing deviation $\DP$ captures temporal information but lacks structural richness. We extend to a 3D state vector representing different aspects of a bounded oscillatory system:

\begin{definition}[S-Entropy Coordinates]
For a bounded oscillatory system, the \emph{S-entropy state} is a 3-tuple $\Stuple = (\Sk, \St, \Se) \in [0,1]^3$ where:
\begin{itemize}[nosep,leftmargin=1em]
  \item $\Sk$ (kinetic entropy): encodes position and velocity within the system's phase space
  \item $\St$ (thermal entropy): encodes the system's energy distribution relative to available states
  \item $\Se$ (energetic entropy): encodes the fraction of maximum possible stored energy
\end{itemize}
All three coordinates are normalized to $[0,1]$ via appropriate scaling.
\end{definition}

\begin{example}
For an ion in an Orbitrap trap:
\begin{align}
\Sk &= \frac{\text{radial position}}{r_{\max}}, \\
\St &= \frac{\text{oscillation frequency}}{\omega_{\max}}, \\
\Se &= \frac{\text{kinetic energy}}{E_{\max}}.
\end{align}
Each coordinate maps a physical quantity to $[0,1]$ independently.
\end{example}

\subsection{Geometry of S-Entropy Space}

\begin{proposition}[Bounded and Measurable]
The S-entropy space $\calS = [0,1]^3 \subset \mathbb{R}^3$ is:
\begin{enumerate}[nosep,label=(\roman*)]
  \item Compact and connected.
  \item Equipped with the standard Lebesgue measure $\mu$ on $\mathbb{R}^3$.
  \item Metric: the Euclidean distance $d(\Stuple_1, \Stuple_2) = \sqrt{(\Sk_1 - \Sk_2)^2 + (\St_1 - \St_2)^2 + (\Se_1 - \Se_2)^2}$.
\end{enumerate}
\end{proposition}

\begin{definition}[S-Entropy Cell]
A \emph{cell} is a Borel-measurable subset $\calC \subseteq \calS$ with positive Lebesgue measure. A \emph{cell partition} $\calP = \{\calC_1, \ldots, \calC_m\}$ covers $\calS$ with disjoint interiors.
\end{definition}

\section{Temporal Programming Model}
\label{sec:temporal}

\subsection{Programs as Timing Cells and Actions}

\begin{definition}[Temporal Program]
A \emph{temporal program} is a tuple $\Pi = (\OOsc, \calP, \mathcal{A})$ where:
\begin{itemize}[nosep,leftmargin=1em]
  \item $\OOsc$ is a reference oscillator,
  \item $\calP = \{\calC_1, \ldots, \calC_m\}$ is a cell partition of $\calS$,
  \item $\mathcal{A}: \calP \to \mathcal{F}$ maps each cell to a pre-compiled action function.
\end{itemize}
\end{definition}

Execution proceeds as follows:
\begin{enumerate}[nosep]
  \item Measure the current S-entropy state: $\Stuple = (\Sk, \St, \Se)$.
  \item Find the unique cell $\calC_i$ containing $\Stuple$.
  \item Dispatch the action $f_i = \mathcal{A}(\calC_i)$.
  \item Execute $f_i$ (which is pre-compiled, pre-verified, and independent of $\Stuple$'s precise value).
\end{enumerate}

\begin{theorem}[Parser-Free Execution]
In temporal programming, execution dispatch depends only on which cell $\Stuple$ occupies, not on the bit pattern or numerical value of $\Stuple$ itself. Therefore:
\begin{enumerate}[nosep,label=(\roman*)]
  \item No payload parser is invoked.
  \item No buffer operations are performed on untrusted data.
  \item The set of executable actions is fixed at compile time: $\mathcal{A}(\calC_1), \ldots, \mathcal{A}(\calC_m)$.
  \item No runtime path leads to injected code.
\end{enumerate}
\end{theorem}

\subsection{Phase Separation}

\begin{definition}[COMPILE/EXECUTE Phases]
The runtime state machine has two phases:
\begin{itemize}[nosep,leftmargin=1em]
  \item \textbf{COMPILE}: Accumulate timing events until a complete trajectory of length $n$ is assembled.
  \item \textbf{EXECUTE}: Determine the cell, dispatch the action, run to completion.
\end{itemize}
\end{definition}

\begin{theorem}[Mutual Exclusion]
At any point in execution, exactly one phase is active. COMPILE and EXECUTE cannot overlap.
\end{theorem}

\begin{proof}
COMPILE requires an open trajectory ($|\tau| < n$, so cell assignment is undefined). EXECUTE requires a closed trajectory ($|\tau| = n$ and cell is known). These are contradictory preconditions.
\end{proof}

\section{Virtual Substate Framework}
\label{sec:virtual}

\subsection{Off-Shell Decomposition}

The key insight is that a single physical state $\Stuple \in [0,1]^3$ can be decomposed into multiple \emph{virtual components} (which may lie outside $[0,1]^3$) as long as their mean equals $\Stuple$.

\begin{definition}[Virtual Substate Decomposition]
For a physical state $\Stuple = (\Sk, \St, \Se) \in [0,1]^3$ and integer $N \geq 2$, a \emph{virtual substate decomposition} is an ordered $N$-tuple of vectors:
\begin{equation}
  \mathbf{V} = (\Stuple_1, \Stuple_2, \ldots, \Stuple_N) \in \mathbb{R}^{3N}
\end{equation}
satisfying the \emph{mean-recovery constraint}:
\begin{equation}
  \frac{1}{N}\sum_{i=1}^{N} \Stuple_i = \Stuple.
  \label{eq:mean-recovery}
\end{equation}
Individual components $\Stuple_i$ may lie outside $[0,1]^3$.
\end{definition}

\begin{proposition}[Existence and Degrees of Freedom]
For any $\Stuple \in [0,1]^3$ and $N \geq 2$:
\begin{enumerate}[nosep,label=(\roman*)]
  \item There exist infinitely many decompositions satisfying \eqref{eq:mean-recovery}.
  \item The degrees of freedom are $3(N-1)$ (the constraint is 3 equations for $3N$ unknowns).
\end{enumerate}
\end{proposition}

\begin{example}
For $N=2$ and $\Stuple = (0.5, 0.5, 0.5)$:
\begin{align}
\Stuple_1 &= (0.7, 0.3, 0.4) \quad (\text{on-shell})\\
\Stuple_2 &= (0.3, 0.7, 0.6) \quad (\text{on-shell})\\
\text{mean} &= (0.5, 0.5, 0.5) \quad (\text{physical state})
\end{align}
or alternatively:
\begin{align}
\Stuple_1 &= (1.2, -0.1, 0.8) \quad (\text{off-shell})\\
\Stuple_2 &= (-0.2, 1.1, 0.2) \quad (\text{off-shell})\\
\text{mean} &= (0.5, 0.5, 0.5) \quad (\text{physical state})
\end{align}
Both are valid decompositions of the same physical state.
\end{example}

\subsection{The Miracle Principle}

\begin{theorem}[Miracle Principle]
For any physical state $\Stuple \in [0,1]^3$, any $N \geq 2$, and any assignment of $(N-1)$ arbitrary virtual components $\Stuple_1, \ldots, \Stuple_{N-1} \in \mathbb{R}^3$, there exists a unique $\Stuple_N$ such that the mean-recovery constraint holds:
\begin{equation}
  \Stuple_N = N\Stuple - \sum_{i=1}^{N-1} \Stuple_i.
\end{equation}
\end{theorem}

This seemingly obvious fact is profound: \emph{any distribution of local off-shell components is compatible with any global physical state}, provided the mean is correct.

\subsection{Physical and Virtual Substates}

\begin{definition}[On-Shell vs Off-Shell]
A component $\Stuple_i$ is:
\begin{itemize}[nosep,leftmargin=1em]
  \item \emph{On-shell}: if $\Stuple_i \in [0,1]^3$.
  \item \emph{Off-shell}: if $\Stuple_i \notin [0,1]^3$ (at least one coordinate outside $[0,1]$).
\end{itemize}
A decomposition is \emph{physical} if all components are on-shell; otherwise it is \emph{virtual}.
\end{definition}

The terminology comes from quantum field theory, where virtual particles have off-shell 4-momenta.

\section{Dual-Domain Architecture}
\label{sec:dualdom}

\subsection{Time-Domain View}

In the temporal-domain view, the program is a standard temporal program:
\begin{itemize}[nosep,leftmargin=1em]
  \item \textbf{Measurement}: Measure $\Stuple(\text{now})$ from the system.
  \item \textbf{Cell lookup}: Find $\calC_i$ containing $\Stuple(\text{now})$.
  \item \textbf{Action dispatch}: Execute $\mathcal{A}(\calC_i)$.
\end{itemize}

This view is:
\begin{itemize}[nosep,leftmargin=1em]
  \item \textbf{Deterministic}: Same state always dispatches same action.
  \item \textbf{Parser-free}: No untrusted data processing.
  \item \textbf{Secure}: Finite attack surface (only the $m$ cell boundaries are sensitive).
\end{itemize}

\subsection{Frequency-Domain View}

In the frequency-domain view, the same execution is reinterpreted as a collection of virtual subtask computations. For each decomposition $\mathbf{V} = (\Stuple_1, \ldots, \Stuple_N)$ satisfying mean-recovery:

\begin{itemize}[nosep,leftmargin=1em]
  \item \textbf{Subtask $i$}: Compute a function $f_i(\Stuple_i)$ for the $i$-th virtual component.
  \item \textbf{Aggregation}: Combine results via mean-recovery: $\text{result} = \mean(f_1(\Stuple_1), \ldots, f_N(\Stuple_N))$.
\end{itemize}

Even though $\Stuple_i$ may be off-shell, the aggregated result is always a valid physical observable (because the mean of any set of values is finite and well-defined).

\begin{theorem}[Dual-Domain Equivalence]
For any temporal program $\Pi$ and any decomposition $\mathbf{V}$ satisfying mean-recovery:
\begin{enumerate}[nosep,label=(\roman*)]
  \item Executing $\Pi$ on $\Stuple$ (time-domain view) produces action $\mathcal{A}(\cell(\Stuple))$.
  \item Decomposing $\Stuple$ into $\mathbf{V}$, executing each subtask, and aggregating (frequency-domain view) produces the same action set.
  \item Both are equivalent from the perspective of the global system.
\end{enumerate}
\end{theorem}

\begin{proof}
By the Miracle Principle, for any decomposition $\mathbf{V}$, $\mean(\mathbf{V}) = \Stuple$ always. The cell $\calC_i$ containing $\Stuple$ is unique. The action $\mathcal{A}(\calC_i)$ depends only on which cell contains the mean, not on the individual components. Therefore, any decomposition produces the same action.
\end{proof}

\subsection{Coordinate Transformation}

The dual-domain architecture is mediated by a Fourier-like coordinate transformation. In the temporal domain, we work with $\Stuple \in \mathcal{S} = [0,1]^3$. In the frequency domain, we work with decompositions $\mathbf{V} \in \mathbb{R}^{3N}$.

The transformation is:
\begin{align}
\text{Time-to-Frequency}: && \Stuple &\mapsto \{\text{all valid decompositions } \mathbf{V}\}\\
\text{Frequency-to-Time}: && \mathbf{V} &\mapsto \mean(\mathbf{V}) = \Stuple
\end{align}

Crucially, the inverse is not unique: many decompositions map to the same $\Stuple$. This redundancy is the source of computational power.

\section{S-Entropy Bidirectional Dijkstra Algorithm}
\label{sec:algorithm}

\subsection{Problem Formulation}

\begin{definition}[Shortest Path in S-Entropy Space]
Given:
\begin{itemize}[nosep,leftmargin=1em]
  \item Current state $\Stuple_{\text{start}} \in [0,1]^3$,
  \item Goal state $\Stuple_{\text{goal}} \in [0,1]^3$,
  \item A temporal program $\Pi$ that generates subtask functions,
\end{itemize}
Find a \emph{path} (sequence of states $\Stuple_0, \Stuple_1, \ldots, \Stuple_k$) that:
\begin{enumerate}[nosep,label=(\roman*)]
  \item Starts at $\Stuple_{\text{start}}$ and ends at $\Stuple_{\text{goal}}$.
  \item Minimizes the total cost: $\sum_{j=0}^{k-1} d(\Stuple_j, \Stuple_{j+1})$.
  \item Respects the temporal program's cell structure (transitions must be valid).
\end{enumerate}
\end{definition}

\subsection{Bidirectional Search Strategy}

Standard Dijkstra searches forward from the start. We introduce a bidirectional approach:

\begin{itemize}[nosep,leftmargin=1em]
  \item \textbf{Forward frontier} (time-domain): Explore reachable states from $\Stuple_{\text{start}}$ via valid transitions.
  \item \textbf{Backward frontier} (frequency-domain): From $\Stuple_{\text{goal}}$, unfold via virtual decompositions to find predecessor states.
  \item \textbf{Meeting point}: Where the two frontiers intersect, the optimal path is found.
\end{itemize}

\subsection{Forward Search (Time-Domain Subtasks)}

\begin{definition}[Valid Transition]
A transition from $\Stuple$ to $\Stuple'$ is \emph{valid} in the temporal program $\Pi$ if:
\begin{itemize}[nosep,leftmargin=1em]
  \item The action $\mathcal{A}(\cell(\Stuple))$ can produce $\Stuple'$ as a possible output state.
  \item The distance $d(\Stuple, \Stuple') \leq \Delta_{\max}$ (bounded by the system's capacitance).
\end{itemize}
\end{definition}

Forward search explores the reachability cone from $\Stuple_{\text{start}}$:
\begin{equation}
  \Reachable_f = \{\Stuple' : \text{valid transition from } \Stuple_{\text{start}} \text{ to } \Stuple'\}.
\end{equation}

\subsection{Backward Search (Frequency-Domain Virtual Substates)}

\begin{definition}[Backward Reachability]
From $\Stuple_{\text{goal}}$, the backward frontier consists of all virtual decompositions $\mathbf{V} = (\Stuple_1, \ldots, \Stuple_N)$ such that:
\begin{enumerate}[nosep,label=(\roman*)]
  \item $\mean(\mathbf{V}) = \Stuple_{\text{goal}}$ (mean-recovery constraint).
  \item At least one component $\Stuple_i$ lies within the reachable cone from $\Stuple_{\text{start}}$ (or in the forward frontier).
\end{enumerate}
\end{definition}

For each decomposition, we extract the ``predecessor directions'': the set of states from which $\Stuple_{\text{goal}}$ could have arisen.

\begin{theorem}[Backward Unfolding]
For any $\Stuple_{\text{goal}} \in [0,1]^3$ and decomposition $\mathbf{V}$ with $\mean(\mathbf{V}) = \Stuple_{\text{goal}}$, the maximum distance from any component to the goal is bounded by $\sqrt{3} \cdot d_{\max}$ where $d_{\max}$ is the diameter of $\calS$.
\end{theorem}

\subsection{The SEBD Algorithm}

\begin{algorithm}
\caption{S-Entropy Bidirectional Dijkstra (SEBD)}
\label{alg:sebd}
\begin{algorithmic}[1]
\Require $\Stuple_{\text{start}}, \Stuple_{\text{goal}} \in [0,1]^3$; Program $\Pi$
\Ensure Path from $\Stuple_{\text{start}}$ to $\Stuple_{\text{goal}}$

\State $Q_f \gets \text{PriorityQueue()}$ \Comment{Forward frontier}
\State $Q_b \gets \text{PriorityQueue()}$ \Comment{Backward frontier}
\State $\text{visited}_f \gets \emptyset$ \Comment{Forward visited}
\State $\text{visited}_b \gets \emptyset$ \Comment{Backward visited}
\State $\text{best\_cost} \gets \infty$
\State $\text{meeting\_point} \gets \text{nil}$

\State $Q_f.\text{push}(\Stuple_{\text{start}}, 0)$
\State $Q_b.\text{push}(\Stuple_{\text{goal}}, 0)$

\While{$Q_f \neq \emptyset$ \textbf{or} $Q_b \neq \emptyset$}
  \If{$Q_f \neq \emptyset$}
    \State $(\Stuple, \text{cost}_f) \gets Q_f.\text{pop}()$
    \If{$\Stuple \in \text{visited}_f$} \textbf{continue} \EndIf
    \State $\text{visited}_f.\text{add}(\Stuple)$

    \If{$\Stuple \in \text{visited}_b$}
      \State $\text{total\_cost} \gets \text{cost}_f + \text{visited}_b[\Stuple]$
      \If{$\text{total\_cost} < \text{best\_cost}$}
        \State $\text{best\_cost} \gets \text{total\_cost}$
        \State $\text{meeting\_point} \gets \Stuple$
      \EndIf
    \EndIf

    \For{each valid transition from $\Stuple$ to $\Stuple'$ with cost $c$}
      \If{$\Stuple' \notin \text{visited}_f$}
        \State $Q_f.\text{push}(\Stuple', \text{cost}_f + c)$
      \EndIf
    \EndFor
  \EndIf

  \If{$Q_b \neq \emptyset$}
    \State $(\mathbf{V}, \text{cost}_b) \gets Q_b.\text{pop}()$ \Comment{$\mathbf{V}$ is decomposition}
    \State $\Stuple_b \gets \mean(\mathbf{V})$
    \If{$\Stuple_b \in \text{visited}_b$} \textbf{continue} \EndIf
    \State $\text{visited}_b.\text{add}(\Stuple_b)$

    \If{$\Stuple_b \in \text{visited}_f$}
      \State $\text{total\_cost} \gets \text{visited}_f[\Stuple_b] + \text{cost}_b$
      \If{$\text{total\_cost} < \text{best\_cost}$}
        \State $\text{best\_cost} \gets \text{total\_cost}$
        \State $\text{meeting\_point} \gets \Stuple_b$
      \EndIf
    \EndIf

    \For{each virtual decomposition $\mathbf{V}'$ from $\mathbf{V}$ with cost $c$}
      \State $Q_b.\text{push}(\mathbf{V}', \text{cost}_b + c)$
    \EndFor
  \EndIf
\EndWhile

\Return Reconstruct path via $\text{meeting\_point}$
\end{algorithmic}
\end{algorithm}

\subsection{Correctness and Termination}

\begin{theorem}[Termination]
Algorithm \ref{alg:sebd} terminates in finite time.
\end{theorem}

\begin{proof}
The state space is $\calS = [0,1]^3$, which is compact. The forward frontier explores reachable states via valid transitions, each with finite cost. The backward frontier explores virtual decompositions, which are discrete (bounded by the temporal program structure). Both frontiers expand monotonically. Since the space is compact and costs are finite, both frontiers eventually reach a meeting point or exhaust the reachable space. The algorithm terminates.
\end{proof}

\begin{theorem}[Optimality]
If Algorithm \ref{alg:sebd} finds a meeting point, the path through that point is optimal (minimum cost).
\end{theorem}

\begin{proof}
The algorithm maintains the invariant that all states in $\text{visited}_f$ have been explored with their true minimum cost from $\Stuple_{\text{start}}$, and all states in $\text{visited}_b$ have been explored with their true minimum cost to $\Stuple_{\text{goal}}$. When a state $\Stuple_m$ is found in both visited sets, the cost $\text{visited}_f[\Stuple_m] + \text{visited}_b[\Stuple_m]$ is the true shortest path cost through that point. Since both frontiers expand greedily by minimum cost, the first meeting point found is globally optimal.
\end{proof}

\begin{theorem}[Mean-Recovery Preservation]
Throughout Algorithm \ref{alg:sebd}, all virtual decompositions $\mathbf{V}$ satisfy the mean-recovery constraint: $\mean(\mathbf{V}) = \Stuple_{\text{goal}}$ (for backward) and transitions in forward search preserve the bounded state space.
\end{theorem}

\section{Security Properties}
\label{sec:security}

\subsection{Temporal-Domain Security}

\begin{theorem}[Structural Incorruptibility]
In the temporal-domain execution of a program $\Pi$, no sequence of S-entropy measurements can cause an action outside $\{\mathcal{A}(\calC_1), \ldots, \mathcal{A}(\calC_m)\}$ to execute.
\end{theorem}

\begin{proof}
Execution is determined entirely by which cell contains $\Stuple$. The cell partition is finite and fixed at compile time. No action outside the predefined set exists for dispatch. Therefore, no measurement can trigger an undefined action.
\end{proof}

\subsection{Replay Attack Resistance}

\begin{theorem}[Monotone Cycle Count Defense]
If a trajectory with timing deviations $(\DP_1, \ldots, \DP_n)$ is recorded at oscillator cycle count $M_0$, and the same bit pattern is replayed at cycle count $M_1 > M_0$, the replayed deviations become $(\DP_1 + \delta, \ldots, \DP_n + \delta)$ where $\delta = (M_1 - M_0)/\fref > 0$. If cell widths are smaller than $\delta$, the replayed trajectory falls in different cells and triggers different actions.
\end{theorem}

In the S-entropy view, this means the replayed state $\Stuple_{\text{replay}}$ differs from the original $\Stuple_{\text{original}}$ by the accumulated phase shift.

\subsection{Computational Transparency}

The frequency-domain view (virtual substates) is mathematically transparent:
\begin{itemize}[nosep,leftmargin=1em]
  \item All decompositions are public (derived mathematically, not hidden).
  \item The mean-recovery constraint is verifiable by any observer.
  \item No cryptographic secrets are needed; security comes from structural properties.
\end{itemize}

\section{Applications}
\label{sec:applications}

\subsection{Real-Time Control Systems}

In industrial control (nuclear plants, drones, manufacturing), the dual-domain model enables:
\begin{itemize}[nosep,leftmargin=1em]
  \item \textbf{Temporal guarantees}: Actions execute deterministically based on cell membership.
  \item \textbf{Rich monitoring}: Virtual decompositions provide early warning signals (off-shell components indicate deviation from nominal).
  \item \textbf{Optimal path planning}: SEBD computes the most efficient sequence of state transitions.
\end{itemize}

\subsection{Distributed Positioning}

In multi-agent positioning (swarms, IoT), each agent:
\begin{enumerate}[nosep]
  \item Maintains its local S-entropy state $(\Sk, \St, \Se)$.
  \item Broadcasts timing pulses (not state values, maintaining privacy).
  \item Uses SEBD to compute shortest path to consensus position.
  \item Synchronization emerges from mean-recovery constraint across the network.
\end{enumerate}

\subsection{Privacy-Preserving Computation}

By decomposing sensitive data into virtual substates (which remain off-shell), computation can proceed without revealing the physical state:
\begin{itemize}[nosep,leftmargin=1em]
  \item Individual virtual components can be held by different agents.
  \item Only the mean-recovery result is observable.
  \item The physical state is never fully reconstructed until final aggregation.
\end{itemize}

\section{Discussion and Open Problems}
\label{sec:discussion}

\subsection{Computational Complexity}

Algorithm \ref{alg:sebd} has complexity:
\begin{itemize}[nosep,leftmargin=1em]
  \item \textbf{Forward frontier}: $O(m \cdot \log V)$ where $m$ is number of cells and $V$ is visited states.
  \item \textbf{Backward frontier}: $O(N \cdot 2^N \cdot \log V)$ where $N$ is decomposition size (exponential in $N$).
  \item \textbf{Overall}: Better than unidirectional for typical cases, but backward complexity grows with $N$.
\end{itemize}

A practical optimization is to limit $N$ (decomposition size) to small values (e.g., $N = 2$ or $N = 4$), making the backward frontier tractable.

\subsection{Extensions}

\begin{itemize}[nosep,leftmargin=1em]
  \item \textbf{Weighted cells}: Assign costs to cell transitions (not uniform cost).
  \item \textbf{Probabilistic transitions}: Allow stochastic jumps between states.
  \item \textbf{Adaptive decomposition}: Choose $N$ dynamically based on search progress.
  \item \textbf{Parallel SEBD}: Run forward and backward frontiers on separate processors.
\end{itemize}

\subsection{Limitations}

\begin{itemize}[nosep,leftmargin=1em]
  \item The algorithm assumes a finite, discrete temporal program. Continuous systems require discretization.
  \item Mean-recovery is a strict constraint; systems violating it are outside the model.
  \item Off-shell components are mathematical abstractions; they have no direct physical realization.
\end{itemize}

\section{Conclusion}
\label{sec:conclusion}

We have presented a dual-domain execution model that unifies temporal programming (parser-free security, deterministic execution) with frequency-domain virtual substates (rich computational structure, off-shell degrees of freedom). The S-Entropy Bidirectional Dijkstra algorithm efficiently finds optimal paths in the 3D S-entropy state space by searching forward in time-domain subtasks and backward through virtual decompositions, meeting at the minimum-cost intermediate state.

Key contributions:
\begin{enumerate}[nosep]
  \item A formal framework for dual-domain computation with mean-recovery constraint.
  \item A novel shortest-path algorithm exploiting bidirectional search.
  \item Proofs of security, correctness, and optimality.
  \item Applications to real-time control, distributed positioning, and privacy-preserving computation.
\end{enumerate}

The model bridges two seemingly incompatible objectives: maximal security (no parser) and rich computation (multi-dimensional state space). Future work includes implementation in real systems, empirical performance evaluation, and extensions to non-Euclidean state spaces.

\begin{thebibliography}{99}

\bibitem{dijkstra1959} E. W. Dijkstra, ``A note on two problems in connexion with graphs,'' \emph{Numer. Math.}, vol. 1, no. 1, pp. 269–271, 1959.

\bibitem{goldberg2005} A. V. Goldberg and C. Harrelson, ``Computing the shortest path: A* search meets graph theory,'' in \emph{Proc. 16th ACM-SIAM Symp. Discrete Algorithms}, 2005, pp. 156–165.

\bibitem{hart1968} P. E. Hart, N. J. Nilsson, and B. Raphael, ``A formal basis for the heuristic determination of minimum cost paths,'' \emph{IEEE Trans. Syst. Sci. Cybern.}, vol. 4, no. 2, pp. 100–107, 1968.

\bibitem{shannon1948} C. E. Shannon, ``A mathematical theory of communication,'' \emph{Bell Syst. Tech. J.}, vol. 27, no. 3, pp. 379–423, 1948.

\bibitem{pontrjagin1962} L. S. Pontryagin, V. G. Boltyanskii, R. V. Gamkrelidze, and E. F. Mishchenko, \emph{The Mathematical Theory of Optimal Processes}. Wiley, 1962.

\bibitem{bellman1957} R. E. Bellman, \emph{Dynamic Programming}. Princeton Univ. Press, 1957.

\bibitem{bhatnagar2009} S. Bhatnagar, R. S. Sutton, M. Ghavamzadeh, and M. Lee, ``Natural actor-critic algorithms,'' \emph{Automatica}, vol. 45, no. 11, pp. 2471–2482, 2009.

\bibitem{konda2000} V. R. Konda and J. N. Tsitsiklis, ``Actor-critic algorithms,'' in \emph{Proc. Adv. Neural Inf. Process. Syst.}, 2000, pp. 1008–1014.

\bibitem{puterman1994} M. L. Puterman, \emph{Markov Decision Processes: Discrete Stochastic Dynamic Programming}. Wiley, 1994.

\bibitem{bertsekas1987} D. P. Bertsekas, \emph{Dynamic Programming: Deterministic and Stochastic Models}. Prentice Hall, 1987.

\bibitem{chvátal1983} V. Chvátal, \emph{Linear Programming}. W.H. Freeman, 1983.

\bibitem{boyd2004} S. Boyd and L. Vandenberghe, \emph{Convex Optimization}. Cambridge Univ. Press, 2004.

\bibitem{nesterov2013} Y. Nesterov, \emph{Introductory Lectures on Convex Optimization: A Basic Course}. Springer, 2013.

\bibitem{calafiore2014} G. C. Calafiore and L. El Ghaoui, \emph{Optimization Models and Applications}. Cambridge Univ. Press, 2014.

\bibitem{rudin1976} W. Rudin, \emph{Principles of Mathematical Analysis}. McGraw-Hill, 3rd ed., 1976.

\bibitem{folland1999} G. B. Folland, \emph{Real Analysis: Modern Techniques and Their Applications}. Wiley, 2nd ed., 1999.

\bibitem{royden1988} H. L. Royden, \emph{Real Analysis}. Macmillan, 3rd ed., 1988.

\bibitem{friedberg2003} S. H. Friedberg, A. J. Insel, and L. E. Spence, \emph{Linear Algebra}. Prentice Hall, 4th ed., 2003.

\end{thebibliography}

\end{document}
